\documentclass[12pt]{article}

\usepackage{amsmath, amssymb}
\usepackage[a4paper, margin=1in]{geometry}
\usepackage{setspace}
\usepackage{hyperref}

\setstretch{1.2}

\begin{document}

\title{Communication is an Equivalence Relation}
\author{Henry Rochester\\
\small Based on lectures by Dr.\ Etienne Pienaar\\
\small Department of Statistical Sciences, University of Cape Town}
\date{}
\maketitle

\section*{Communication is an Equivalence Relation}

We show that the communication relation between states of a Markov chain is an equivalence relation. In particular, we show the transitivity property. C'est parti!\\

\textbf{Statement:}

If $i \leftrightarrow j$ and $i \leftrightarrow k$, then $j \leftrightarrow k$.

\vspace{0.5cm}

\subsection*{Proof (exam-style version)}

Suppose that $i \leftrightarrow j$ and $i \leftrightarrow k$.\\

By definition of communication, there exist integers $v,w \geq 1$ such that:\\
\[
\gamma_{ij}(v) > 0
\qquad \text{and} \qquad
\gamma_{ji}(w) > 0.
\]\\

Similarly, there exist integers $x,y \geq 1$ such that:\\
\[
\gamma_{ik}(x) > 0
\qquad \text{and} \qquad
\gamma_{ki}(y) > 0.
\]\\

To show $j \to k$, consider the transition probability $\gamma_{jk}(w+x)$. By the Chapman-Kolmogorov equations,\\
\[
\gamma_{jk}(w+x)
=
\sum_{r \in \cal{U}} \gamma_{jr}(w)\gamma_{rk}(x).
\]\\

Since $i \in \cal{U}$, one of the terms in the sum is given by: \\
\[
\gamma_{ji}(w)\gamma_{ik}(x).
\]\\

Both factors are strictly positive, thus it follows that: \\
\[
\gamma_{ji}(w)\gamma_{ik}(x) > 0.
\]\\

Hence,\\
\[
\gamma_{jk}(w+x)
\geq
\gamma_{ji}(w)\gamma_{ik}(x)
>0.
\]\\

Thus,\\
\[
\gamma_{jk}(w+x)>0,
\]\\
and therefore $j \to k$.\\


Similarly, to show $k \to j$, consider $\gamma_{kj}(y+v)$. By Chapman-Kolmogorov,\\
\[
\gamma_{kj}(y+v)
=
\sum_{r\in \cal{U}}\gamma_{kr}(y)\gamma_{rj}(v).
\]\\

Again, the term corresponding to $r=i$ is:
\[
\gamma_{ki}(y)\gamma_{ij}(v).
\]\\

Both factors are strictly positive, so:
\[
\gamma_{ki}(y)\gamma_{ij}(v)>0.
\]\\

Thus,\\
\[
\gamma_{kj}(y+v)
\geq
\gamma_{ki}(y)\gamma_{ij}(v)
>0.
\]\\

Hence,\\
\[
\gamma_{kj}(y+v)>0,
\]\\
which implies $k \to j$.\\

Since $j \to k$ and $k \to j$, we conclude that: \\
\[
j \leftrightarrow k.
\]\\

\[
\boxed{j \leftrightarrow k}
\]

as required. \hfill $\square$

\newpage

\section*{Proof with Intuition and Detailed Reasoning}

The goal of this proof is to show that the communication relation between states behaves like an equivalence relation. In particular, we want to prove \textit{transitivity}. That is to say:\\

\begin{quote}
If state $i$ communicates with state $j$, and state $i$ communicates with state $k$, then states $j$ and $k$ must also communicate with each other.\\
\end{quote}

\vspace{0.3cm}

We begin by recalling the definition of communication. Saying that $i \leftrightarrow j$ means two things:\\

\begin{itemize}
\item it is possible to go from state $i$ to state $j$ in some number of steps, and
\item it is possible to go from state $j$ back to state $i$ in some number of steps.
\end{itemize}

Formally, this means that there exist integers $v,w \geq 1$ such that:\\
\[
\gamma_{ij}(v)>0
\qquad \text{and} \qquad
\gamma_{ji}(w)>0.
\]\\

Intuitively, this simply means that there exists a path from $i$ to $j$ and another path from $j$ back to $i$.\\

\vspace{0.3cm}

Similarly, the assumption $i \leftrightarrow k$ tells us that there exist integers $x,y \geq 1$ such that:\\
\[
\gamma_{ik}(x)>0
\qquad \text{and} \qquad
\gamma_{ki}(y)>0.
\]\\

So we know that we can travel between $i$ and $k$ in both directions as well.\\

\vspace{0.3cm}

Now we want to show that $j$ can reach $k$. Intuitively, this should be true because we already know the following:

\begin{enumerate}
\item we can go from $j$ to $i$ in $w$ steps, and
\item we can go from $i$ to $k$ in $x$ steps.
\end{enumerate}

Therefore, by simply concatenating these two journeys, we obtain a path from $j$ to $k$ that takes $w+x$ steps.\\

To formalise this idea, we use the Chapman-Kolmogorov equations. These allow us to decompose the probability of moving from one state to another over multiple steps by summing over intermediate states. Specifically,\\
\[
\gamma_{jk}(w+x)
=
\sum_{r\in \cal{U}}\gamma_{jr}(w)\gamma_{rk}(x).
\]\\

This equation tells us that the probability of moving from $j$ to $k$ in $w+x$ steps can be computed by summing over all intermediate states $r$ that the chain might occupy after $w$ steps.\\

Since $i$ is one of the states in the state space $U$, one of the terms in this sum is:\\
\[
\gamma_{ji}(w)\gamma_{ik}(x).
\]\\

But we know that both of these probabilities are strictly positive. Therefore their product is also strictly positive. Consequently, the entire sum must also be positive, which implies\\
\[
\gamma_{jk}(w+x)>0.
\]\\

Thus $j \to k$.\\

\vspace{0.3cm}

A completely symmetric argument shows that $k \to j$. We simply reverse the roles of the paths:\\

\begin{itemize}
\item first move from $k$ to $i$ in $y$ steps,
\item then move from $i$ to $j$ in $v$ steps.
\end{itemize}

Using Chapman-Kolmogorov again gives: \\
\[
\gamma_{kj}(y+v)
=
\sum_{r\in \cal{U}}\gamma_{kr}(y)\gamma_{rj}(v),
\]\\
and the term corresponding to $r=i$ is strictly positive. Hence: \\
\[
\gamma_{kj}(y+v)>0,
\]\\
so $k \to j$.\\

\vspace{0.3cm}

Putting these two results together yields: \\
\[
j \to k
\qquad \text{and} \qquad
k \to j.
\]\\

Therefore,\\
\[
j \leftrightarrow k.
\]

This proves that the communication relation is transitive.

\section*{Final Intuition}

The entire proof is essentially a statement about paths in the Markov chain:\\

\begin{quote}
If $j$ can reach $i$, and $i$ can reach $k$, then by following those paths one after the other we obtain a path from $j$ to $k$.
\end{quote}

The Chapman-Kolmogorov equations simply provide the formal mathematical justification that the probability of such a combined path must be positive.

\section*{References}

\begin{itemize}

\item Pienaar, E. (University of Cape Town). 
\textit{Stochastic Processes}. 
Lecture notes and in-person lectures, Department of Statistical Sciences,
University of Cape Town.

\end{itemize}

\end{document}